\documentclass[12pt]{article}
\usepackage{fontspec}
\usepackage{xunicode}
\usepackage{polyglossia}
\setmainlanguage{french}

\usepackage{reledmac}%Faire une édition critique
\usepackage{reledpar}%Mettre deux textes en vis-à-vis

\begin{document}

\begin{pairs}



\begin{Leftside}
	\beginnumbering

	 \pstart	 

Rollant est proz e Oliver est sage : \edtext{Ambedui}{\Afootnote{entrambi \textit{ms.} de Venise}}
ambedui unt merveillus vasselage ; Puis que il sunt a cheval e as armes, ja pur murir n'eschiverunt bataille.

\pend
\pstart
Bon sunt li cunte e lur paroles haltes. Felun paien par grant irur chevalchent.
\edtext{Dist Oliver : « Rollant, veez en alques !}{\Afootnote{{omit.} \textit{ms.} de Châteauroux}}
Dist Oliver : « Rollant, veez en alques ! Cist nus sunt pres, mais trop nus est loinz Carles. Soner tun olifan, si l'orrat li reis. »
\pend
\pstart

Rollant respunt : « Jo fereie que fols ! En dulce France en perdreie mun los. Sempres ferrai de Durendal granz colps.
Sempres ferrai de Durendal granz colps : sanglant en ert li branz entresqu'a l'or. Felun paien mar i vindrent as porz : jo vos plevis, tuz sunt jugez a mort. »
\pend
\endnumbering
\end{Leftside}

%Ici se trouve la traduction moderne (prose, d'après J. Dufournet, 1993)

\begin{Rightside}
\beginnumbering

\pstart


Roland est vaillant et Olivier est sage : tous deux font preuve d'une bravoure admirable ; dès lors qu'ils sont à cheval et en armes, jamais, dussent-ils en mourir, ils ne fuiront la bataille.

\pend

\pstart
Les deux comtes sont excellents, et leurs paroles sont fières. Les païens félons chevauchent avec une grande fureur. Olivier dit : « Roland, regarde un peu ! Ils sont proches de nous, mais Charles est trop loin. Sonne ton olifant, et le roi t'entendra. »
\pend

\pstart
Roland répond : « Ce serait agir en insensé ! En douce France j'en perdrais ma renommée. À l'instant je frapperai de Durendal de grands coups : la lame en sera ensanglantée jusqu'à l'or. Les félons païens ont eu tort de venir aux défilés : je vous le jure, ils sont tous condamnés à mort. »
\pend

\endnumbering

\end{Rightside}

\end{pairs}

    \Columns

\end{document}